\section{Design}
\label{sec:design}

\subsection{Policy Model}

ActPlane policies are object-level information-flow rules over three node
types: processes, files, and endpoints. A process node represents a Linux task
identified by pid and command identity. A file node is currently keyed by path.
An endpoint node is currently keyed by IPv4 address. The model is intentionally
coarser than byte-level dynamic taint tracking~\cite{kemerlis2012libdft}; it
trades precision for low-overhead whole-system coverage.

\subsubsection{Labels and Propagation}

Each node carries a bitset of labels. A source declaration introduces labels:
an exec source labels processes that execute a matching binary, a file source
labels matching files, and an endpoint source labels matching network
endpoints. Labels propagate through fixed transfer functions:

\begin{itemize}
\item fork: a child process inherits the parent's labels;
\item exec: a process absorbs labels from the executed file and matching exec
sources;
\item read: a process absorbs labels from the file it reads;
\item write: a file absorbs labels from the process writing it;
\item connect: an endpoint records labels from the connecting process.
\end{itemize}

This is the provenance invariant ActPlane maintains: if a process carries label
\texttt{SECRET}, then information from some \texttt{SECRET} source has reached
that process through observed OS-level flow edges.

\subsubsection{Rules, Gates, and Effects}

A rule contains one or more deny clauses. A clause specifies an operation
(\texttt{exec}, \texttt{read}, \texttt{write}, \texttt{unlink},
\texttt{connect}, or \texttt{open}), a target pattern, a Boolean expression over
labels, and an optional condition such as a target allow-list, a lineage gate,
or an after gate.

\begin{figure}[t]
\begin{verbatim}
label AGENT
source SECRET = file "**/.env"

rule no-secret-egress:
  deny connect endpoint "*" if SECRET
  effect block
  reason "secret-derived data must stay local"
  remediation "run the redactor or ask the user"

rule no-agent-git:
  deny exec "**/git" if AGENT
  effect kill
  reason "agent runs must not invoke git"
\end{verbatim}
\caption{A shortened ActPlane policy. The first rule asks for LSM-only
pre-operation deny. The second rule asks for harness-level termination and can
fire from tracepoint observation.}
\Description{A code listing showing two ActPlane rules: one blocking secret
egress and one killing git execution by an agent.}
\label{fig:policy-example}
\end{figure}

Rules have explicit effects.

\textbf{Audit.} The kernel reports the violation and allows the action to
proceed. This is useful for soft workflow guidance and for measurement.

\textbf{Block.} A BPF-LSM hook returns \texttt{-EPERM} before the operation is
committed. \texttt{block} is not supported in tracepoint mode and is not
silently downgraded to \texttt{audit} or \texttt{kill}.

\textbf{Kill.} The kernel sends \texttt{SIGKILL} to the current task. If a
pre-operation LSM hook is active, ActPlane can also return \texttt{-EPERM};
from a tracepoint, \texttt{kill} is harness-level termination rather than
pre-operation denial.

If several supported rules match the same event, ActPlane chooses the strongest
effect in the order \texttt{kill > block > audit}. Unsupported effects are
ignored by that backend, so a tracepoint \texttt{audit} rule can still fire even
when an unsupported \texttt{block} rule also matches.

\subsubsection{Why This Is a Harness DSL}

The language is deliberately not a general-purpose MAC policy. Its primitives
map to agent workflows: ``after tests,'' ``through a review tool,'' ``data
derived from untrusted input,'' and ``read-only sub-agent.'' The grammar is
small because policy authors should be able to translate natural-language
instructions into a few recurring source/sink/gate shapes.

\subsection{System Architecture}

Figure~\ref{fig:architecture} summarizes ActPlane's control flow. The Rust
runner discovers a project policy, compiles its embedded DSL into a fixed binary
configuration, starts the eBPF loader, seeds the target command's pid lineage
with the \texttt{AGENT} label, and streams violation events back into a feedback
file.

\begin{figure}[t]
\begin{verbatim}
actplane.yaml
  policy: |
        |
        v
Rust compiler ------> struct taint_config
        |                    |
        v                    v
 actplane run -----> eBPF loader + BPF maps
        |                    |
        v                    v
  target command <--- kernel hooks
        |
        v
.actplane/last-violation.txt
\end{verbatim}
\caption{ActPlane pipeline. The collector compiles policy metadata and reasons
in userspace, while the kernel enforces and reports rule identifiers.}
\Description{ASCII diagram of policy compilation, eBPF loading, target command
execution, and feedback output.}
\label{fig:architecture}
\end{figure}

\subsubsection{Compiler and Kernel ABI}

The collector parses the DSL and lowers it into a C-compatible
\texttt{struct taint\_config}. Lowering allocates label bits, expands Boolean
expressions into required and forbidden masks, lowers glob patterns into exact,
prefix, suffix, or any match kinds, and compiles IPv4 endpoint patterns into
network/mask pairs. Human-readable rule names, reasons, and remediation strings
remain in Rust metadata; the kernel only receives fixed-size rule tables and
emits rule identifiers.

This split keeps the eBPF program verifier-friendly. The kernel program performs
bounded loops over fixed-size arrays in read-only data, copies rule entries into
stack locals before matching, and uses explicit bounds checks for argv scanning
and pattern matching.

\subsubsection{Kernel State}

The kernel backend maintains three principal maps:

\begin{itemize}
\item \texttt{ts\_proc}: pid to labels and lineage gates;
\item \texttt{ts\_root}/\texttt{ts\_sess}: root process and session-level after
gates;
\item \texttt{ts\_file} and \texttt{ts\_endp}: file and endpoint labels.
\end{itemize}

Tracepoints maintain lineage and perform observation-compatible checks. BPF-LSM
programs perform pre-operation checks for exec, file access/mutation, and socket
connect. When BPF-LSM is not active in \texttt{/sys/kernel/security/lsm}, the
loader disables LSM programs and warns that \texttt{block} rules will not fire.

\subsubsection{Unified Event Handling}

Each hook normalizes kernel arguments into a small event structure: subject pid,
object kind, access kind, backend mode, target display string, optional argv
blob, and optional IPv4 address. The shared handler resolves current labels,
adds object-source labels for read/connect events, evaluates supported rules,
emits exactly one \texttt{TAINT\_VIOLATION} event when a supported rule matches,
and then applies propagation updates.

\subsection{Corrective Feedback}

The kernel event contains the rule id, effect, target, pid, and flags indicating
whether the operation was blocked or killed. The Rust runner maps the rule id
back to a rule name, reason, operation, and remediation string, then appends a
model-facing payload to \texttt{.actplane/last-violation.txt}. The same binary
implements \texttt{actplane feedback-hook}, a post-tool hook adapter that reads
new bytes from that file and returns them as additional context to compatible
agents.

The feedback path does not re-evaluate policy in userspace. The kernel is the
authority for rule matches. Userspace only formats the reason so a cooperative
agent can choose a different route.


\section{Implementation}

ActPlane is implemented as a Rust collector and an eBPF loader/program pair.
The current codebase builds a single \texttt{actplane} binary that embeds the
compiled loader, extracts it at runtime, compiles project policy YAML, and runs
a target command under the policy harness.

\subsection{Runner}

The command \texttt{actplane run -- <cmd>} discovers \texttt{actplane.yaml}
upward from the current directory unless \texttt{--policy} or \texttt{--rule} is
specified. The policy YAML contains a \texttt{policy: |} DSL block and an
optional feedback path. Run mode requires an \texttt{AGENT} label; before the
target begins executing, the loader seeds the stopped target pid with that label
so all descendants inherit the agent identity.

The runner prepares \texttt{.actplane/last-violation.txt} and a hook state file,
exports their paths to the child, waits until the loader reports readiness, then
continues the target process. This ensures the first target action is observed.

\subsection{BPF Backends}

The eBPF program attaches to process lifecycle tracepoints for fork, exec, and
exit, to syscall tracepoints for observation-compatible file and connect events,
and to BPF-LSM hooks for pre-operation enforcement. The loader checks whether
the running kernel has the \texttt{bpf} LSM active. If not, it uses tracepoint
mode. In tracepoint mode, supported effects are exactly \texttt{audit} and
\texttt{kill}; \texttt{block} is unsupported because there is no pre-operation
verdict to return.

\subsection{Current Limitations}

The implementation deliberately exposes several limitations rather than hiding
them behind ambiguous semantics.

\textbf{Path identity.} File labels are currently keyed by path hash. This is
sufficient for the prototype but should move toward inode/device identity for a
stronger claim.

\textbf{Endpoint identity.} Kernel connect matching is numeric IPv4. Hostname,
DNS, SNI, and HTTP-layer policies require userspace assist or additional
instrumentation.

\textbf{Exec arguments.} \texttt{@arg} predicates are currently observed after
exec unless an argv cache is added before \texttt{bprm\_check\_security}. Thus
\texttt{effect block} with an argv predicate is not available through the
tracepoint path; \texttt{effect audit} can report and \texttt{effect kill} can
terminate after exec.

\textbf{Precision.} Labels propagate at object granularity. This can over-taint
compared with byte-level DIFT systems. Section~\ref{sec:evaluation} measures
false positives and validates that declassification and gate paths prevent
unnecessary failures.
